\problem{P19}{跳过前段后的最佳选择}{记录与最优停止}
$n\ge2$ 名候选人按均匀随机顺序到达，水平互不相同，只能观察当前候选人在已出现者中的相对排名，拒绝后不能召回。采用规则：先拒绝前 $r$ 人，再选择第一个超过此前所有人的候选人；若此后无人创新高，则选择失败。$0\le r\le n-1$，其中 $r=0$ 表示直接选第一人。求选中全体最佳者的概率，并给出最优 $r$ 的确定方法及大样本极限。
\solution $r=0$ 时成功率为 $1/n$。对 $r\ge1$，若最佳者位于 $k>r$，要一直等到他，前 $k-1$ 人中的最佳者必须在被跳过的前 $r$ 人中。其位置均匀，所以条件成功概率为 $r/(k-1)$。因此
\[
\boxed{P_r=\frac rn\sum_{j=r}^{n-1}\frac1j,\quad1\le r\le n-1.}
\]
相邻两值的差为
\[
P_{r+1}-P_r=\frac1n\left(\sum_{j=r+1}^{n-1}\frac1j-1\right),
\quad1\le r\le n-2.
\]
括号随 $r$ 严格递减，故成功率先升后降；逐次增加 $r$，直到该尾部调和和不超过 $1$ 即可确定峰值，等于 $1$ 时相邻两个值并列。还须比较 $r=0$；因 $P_1=H_{n-1}/n\ge1/n$，它不会严格优于 $r=1$。

若 $r/n\to x\in(0,1)$，调和和趋于 $-\log x$，故 $P_r\to-x\log x$。该函数在 $x=e^{-1}$ 唯一取最大值 $e^{-1}$。结合上述离散差分条件，得到最优跳过比例趋于 $1/e$，最优成功率也趋于 $1/e$。
\insight 条件化于最佳者到达时刻，把整段选择历史压缩为此前最高者的位置。优化阶段只需考察一个单调的相邻差分。
